\documentclass{article}
\usepackage{titling}
\usepackage{ctex}
\usepackage{tikz}
\usepackage{graphicx}
\usepackage{amsmath}  % 修复 \text 报错
\usetikzlibrary{positioning}

\tikzset{
    box/.style={draw, thin, align=center, inner sep=4pt, outer sep=3pt, font=\small},
    boxred/.style={draw=red, thin, align=center, inner sep=4pt, outer sep=3pt, font=\small}
}

\title{数据结构\\ \large 排序}
\author{}
\date{}

\begin{document}
\maketitle

\section{归并排序}
归并排序基于分治思想，先将序列不断分解为子序列，再将有序子序列两两合并，最终得到完整有序序列。

\begin{figure}[htbp]
\centering
\resizebox{0.9\textwidth}{!}{%
\begin{tikzpicture}[
    node distance=0.6cm and 0.2cm,
    >=stealth,
    thick,
    every path/.style={thick}
]

% ================== 顶层 ==================
\node[box] (top) at (0,7) {10 6 7 1 3 9 4 2};

% ================== 第二层 ==================
\node[box] (l2l) at (-4.2,5.8) {10 6 7 1};
\node[box] (l2r) at (4.2,5.8) {3 9 4 2};

\draw (top) -- (l2l);
\draw (top) -- (l2r);

% ================== 第三层 ==================
\node[box] (l3l1) at (-5.6,4.6) {10 6};
\node[box] (l3l2) at (-2.8,4.6) {7 1};
\node[box] (l3r1) at (2.8,4.6) {3 9};
\node[box] (l3r2) at (5.6,4.6) {4 2};

\draw (l2l) -- (l3l1);
\draw (l2l) -- (l3l2);
\draw (l2r) -- (l3r1);
\draw (l2r) -- (l3r2);

% ================== 第四层：单个元素 ==================
\node[box] (l4l1) at (-6.6,3.4) {10};
\node[box] (l4l2) at (-4.6,3.4) {6};
\node[box] (l4l3) at (-3.6,3.4) {7};
\node[box] (l4l4) at (-1.6,3.4) {1};
\node[box] (l4r1) at (1.6,3.4) {3};
\node[box] (l4r2) at (3.6,3.4) {9};
\node[box] (l4r3) at (4.6,3.4) {4};
\node[box] (l4r4) at (6.6,3.4) {2};

\draw (l3l1) -- (l4l1);
\draw (l3l1) -- (l4l2);
\draw (l3l2) -- (l4l3);
\draw (l3l2) -- (l4l4);
\draw (l3r1) -- (l4r1);
\draw (l3r1) -- (l4r2);
\draw (l3r2) -- (l4r3);
\draw (l3r2) -- (l4r4);

% ================== 绿色分割线 ==================
\draw[green, thick] (-7.5,2.7) -- (7.5,2.7);

% ================== 左侧文字 ==================
\node[align=right, font=\small] at (-8.2,3.2) {分解};
\node[align=right, font=\small] at (-8.2,1.6) {合并};

% ================== 第五层：第一次合并 ==================
\node[boxred] (l5l1) at (-5.6,1.8) {6 10};
\node[boxred] (l5l2) at (-2.8,1.8) {1 7};
\node[boxred] (l5r1) at (2.8,1.8) {3 9};
\node[boxred] (l5r2) at (5.6,1.8) {2 4};

\draw (l4l1) -- (l5l1);
\draw (l4l2) -- (l5l1);
\draw (l4l3) -- (l5l2);
\draw (l4l4) -- (l5l2);
\draw (l4r1) -- (l5r1);
\draw (l4r2) -- (l5r1);
\draw (l4r3) -- (l5r2);
\draw (l4r4) -- (l5r2);

% ================== 第六层 ==================
\node[boxred] (l6l) at (-4.2,0.6) {1 6 7 10};
\node[boxred] (l6r) at (4.2,0.6) {2 3 4 9};

\draw (l5l1) -- (l6l);
\draw (l5l2) -- (l6l);
\draw (l5r1) -- (l6r);
\draw (l5r2) -- (l6r);

% ================== 第七层：最终结果 ==================
\node[boxred] (l7) at (0,-0.6) {1 2 3 4 6 7 9 10};

\draw (l6l) -- (l7);
\draw (l6r) -- (l7);

\end{tikzpicture}
}
\caption{归并排序示意图}
\end{figure}

有序子数组合并过程示例（A = [6, 10]，B = [1, 7]）：
\begin{enumerate}
    \item 初始状态：指针 $i$ 指向 A 第一个元素：$i=0,\ A[i]=6$；指针 $j$ 指向 B 第一个元素：$j=0,\ B[j]=1$；结果数组 $\text{res} = []$。
    \item 对比 $6$ 和 $1$，$1$ 更小 $\to$ 把 $1$ 放进 $\text{res}$，$j \mathrel{+}= 1$。$\text{res}$：$[1]$，现在 $B[j]=7$。
    \item 对比 $6$ 和 $7$，$6$ 更小 $\to$ 把 $6$ 放进 $\text{res}$，$i \mathrel{+}= 1$。$\text{res}$：$[1, 6]$，现在 $A[i]=10$。
    \item 对比 $10$ 和 $7$，$7$ 更小 $\to$ 把 $7$ 放进 $\text{res}$，$j \mathrel{+}= 1$；此时 B 数组全部取完（$j$ 走到末尾）。$\text{res}$：$[1, 6, 7]$。
    \item B 数组空了，直接把 A 剩下的元素全部追加进来，剩下 A 里的 $10$。$\text{res}$：$[1, 6, 7, 10]$。
\end{enumerate}

\end{document}
